
% Our study focuses on "Can artifacts in public registries be verified", not "Are artifacts in public registries authentic". The first question focuses on two things: Is it possible to recover the correct source commit, build environment, and build instructions from which a package was built from, and if we rebuild using the recovered information, what level of artifact equivalence do we get?


\section{Knowledge Gaps and Research Questions}
\label{sec:knowledge-gaps}

% Towards Measuring Supply Chain Attacks on Package Managers for Interpreted Languages (2021) - includes attack vectors like typosquatting and publishing but no detail how existing packages are compromised.
% Backstabber’s Knife Collection: A Review of Open Source Software Supply Chain Attacks (2019) - focused on anatomy of attacks (ecosystems affected, how they are triggered, conditional execution, operating system affected, etc)


% \begin{table}
%     \centering
%     \caption{Attack vectors of sampled software supply chain incidents involving existing package compromise. }
%     \label{tab:ssc-incidents}
%     \begin{tabular}{p{2.5cm}|p{2.5cm}p{0.8cm}p{0.8cm}}
%       \toprule
% \textbf{Attack vector} &
% \textbf{Total} &
% \textbf{Src. modified} &
% \textbf{Valid attestation} \\
% \midrule
% Submit bad code
%   & 1
%   & 1
%   & 0 \\

% Compromise source 
%   & 2 
%   & 2
%   & 0 \\

% Compromise build
%   & 4 (NPM: 1; PyPI: 2)
%   & 0
%   & 1 \\

% Bypass CI/CD 
%   & 17 (NPM: 13; PyPI: 4)
%   & 0
%   & 2 \\
% \midrule
% \textbf{Total}
%   & \textbf{24}
%   & \textbf{3}
%   & \textbf{3} \\
% \bottomrule
%     \end{tabular}
% \end{table}


% \subsection{Motivation}
% \label{subsec:motivation}

% \comment[KC]{I understand the job the motivation is trying to achieve. But as a reader I can't help but feel it should come earlier than this or later than this. A. Earleir: It feels like just after we mention that build infra is a attractive target for attackers + define how verifiability and Reproducible builds work-> we can then go ahead and use this motivation as a case study context to say, see in the past few years these attacks could have been prevented this way etc. B. After (if we could make this a key part of our study method design): It also feel like it could come around the methods  maybe to define why we masde some metodological choice like registry selection etc. In this current position it feels disruptive to the flow of the paper }

% \cref{subsec:bg-repro-builds} identify reproducible and verifiable builds as defenses against build compromise attacks.
% However, prior studies~\cite{} do not show what proportion of real-world attacks would have been prevented by these techniques.
% To examine this gap and motivate further studies, we collected public reports of attacks that compromised existing packages between 2021 and 2026 and sampled 24 incidents.
% For each incident, we recorded the affected ecosystem, affected packages and versions, threat vector in \cref{fig:ssc-threats}, whether the upstream source repository was modified, and whether the malicious package had a valid attestation.

% \Cref{tab:ssc-incidents} summarizes our results, with full details in our artifacts.
% We find that 21 of 24 incidents affected interpreted-language ecosystems, including NPM, PyPI, and PHP.
% We also find that 21 of 24 incidents introduced malicious code through the build or publication process without corresponding source changes.
% Thus, artifact verifiability could have detected source-to-artifact mismatches in 87.5\% of incidents, covering more than a thousand compromised artifacts.
% Furthermore, 3 incidents affected packages with attestations, demonstrating that attestations alone are insufficient because attackers can abuse legitimate release paths to publish malicious artifacts with valid attestations.
% Together, these findings motivate artifact verifiability as a stronger mechanism for validating published artifacts.

% Three incidents illustrate this boundary.
% In LiteLLM, attackers used compromised PyPI credentials to publish malicious versions directly to the registry.
% In the TanStack Mini Shai-Hulud incident, attackers compromised the release pipeline and produced malicious packages with valid SLSA provenance.
% Artifact verifiability can detect both attacks by identifying mismatches between the published artifacts and expected source.
% In contrast, the XZ Utils backdoor was introduced into the upstream repository, so artifact verifiability would not detect it.





% \input{sections/supply_chain_ind_table}
% Recent supply-chain incidents demonstrate that trust in released package artifacts cannot be reduced to trust in source repositories alone. 
% Table~\ref{tab:recent-supply-chain-incidents} summarizes representative incidents from the last four years and classifies whether source--artifact verification or provenance enforcement could have helped detect or reject the compromised artifact before consumption.

% The \texttt{event-stream} incident demonstrated that malicious code can enter downstream applications through packages distributed by the npm registry~\cite{npm-event-stream-2018}.
% In the Python ecosystem, PyPI's analysis of the malicious \texttt{aiocpa} package found obfuscated credential-exfiltration code in published distributions but no evidence of the same code in the linked source repository~\cite{pypi-aiocpa-2024}.
% The \texttt{ultralytics} incident further showed how the compromise of a GitHub Actions workflow and publishing credentials could result in malicious PyPI artifacts, while also illustrating how attestations can narrow the scope of an investigation by revealing the path through which a release was published~\cite{pypi-ultralytics-2024}.
% Beyond language-specific registries, the \texttt{xz} backdoor demonstrated that release archives may contain build logic absent from the public repository~\cite{nvd-xz-2024}.

% These incidents expose a fundamental gap between the source code that users and auditors inspect and the artifacts that package registries actually distribute.
% They also illustrate that no single mechanism is sufficient.
% Source--artifact verification can reveal unexplained differences between public source and released artifacts, while provenance can provide evidence about the identity and workflow used to publish a release.
% However, provenance alone does not establish that a package is trustworthy, and a reproducible build is only useful when an independent verifier can identify the relevant source revision and reconstruct the required build conditions.

% This gap motivates our central question:
% \emph{Can artifacts distributed through public package registries be independently traced to their source revisions and rebuilt in recovered environments such that the rebuilt artifacts are equivalent to their published counterparts?}
% \comment[Oreofe]{As Table~\ref{tab:recent-supply-chain-incidents} shows, recent supply-chain attacks frequently exploit weak links between source code, release workflows, and the artifacts distributed through registries. However, the table also shows that reproducible builds are not a complete defense: their effectiveness depends on whether an independent verifier can recover the relevant source revision, build environment, and publishing workflow.}

  % \comment[Oreofe]{Attested and heuristic samples are independent populations (mostly the case), not paired. Our L1 gap confounds provenance value with package-population differences (attested packages skew newer / better tooled). Crates is fine (Trustpub vs vcs\_info are same packages). For PyPI/npm we either need to (a) reframe RQ3 as population comparison and acknowledge possible confound, or (b) run heuristic recovery on the attested sample for a paired comparison. Preferences?}

% \subsection{Knowledge Gap and Research Questions}

Verifiable builds promise to mitigate software supply-chain attacks that exploit compromised build and release pipelines.
However, existing verifiable-build systems either operate in ecosystems with centralized build infrastructure, where rebuild metadata is recorded during build, or rely on human intervention to recover missing metadata in decentralized ecosystems.
As a result, we do not yet know whether artifacts in decentralized-build ecosystems can be verified automatically.

% \JD{One example of many places where you need to prefer the phrase `decentralized-build'}
This paper asks: \textit{To what extent can published artifacts in decentralized-build package ecosystems be independently verified against their source without human intervention?}
We structure this question into four research questions:

\begin{description}
\JD{RQ1 includes the ablation on metadata extraction and whatever other knobs you've got}
\JD{Does combining strategies produce better results, or is one strategy dominant}

\item[\textbf{RQ1.}] \textit{Verifiability results:} What proportion of registry artifacts can be verified against their source?

\item[\textbf{RQ2.}] \textit{Root causes:} What are the root causes of artifact verifiability failures?

\item[\textbf{RQ3.}] \textit{Factor analysis:} How is verifiability affected by artifact characteristics and development choices?

%\item[\textbf{RQ4.}] How do artifact characteristics and build-system choices affect verifiability?

\item[\textbf{RQ4.}] \textit{Comparison:} How does \toolname compare with state-of-practice rebuild systems?
\end{description}


% It seems LastPyMile focused on just measuring the differences between the source and package, and produces results to support verification choices.
% However, we don't want to measure differences. We want to measure what proportion of packages are verifiable and what choices will make them verifiable. Their result is that the code distributed in registries has many changes with respect to the code stored in repositories.
% 15% of python files have differences in init and setup files, which could have been added by the build tools during the packaging processes. Reproducible builds corrects for these cases.

% ICSE 25: (1) They used the default build command to build the package. (2) They discarded packages they could not identify their source repository or where build failures occurred (~800/4000). (3) They use reprotest, which requires fixing the build infrastructure specifications (toolchain, versions, dependencies, operating system) (4) They acknowledged that identifying the source revision for a package is non-trivial and subject of cited ongoing research. Also identified that environment dependnecies (compiler versions, system libraries, operating system configurations) can vary across build environments and cause potential discrepancies. (5) Independent build of debian packages show that only 30% can be reproduced while about 90% can be reproduced using reprotest. (6) Timestamp was the major source of non-reproducibility, which was fixed by modifying the build instructions. (7) They conclude that native code extensions do not impact reproducibility, but this is with respect to their controlled environment.
% \PA{At the start of this section, we need to describe the unique characteristics of source-oriented packages that make tiered equivalence levels feasible. At the end, we should describe how our tiers differ from prior work.}
% \comment[Oreofe]{We also admit line-ending (CRLF vs.\ LF) differences at this level via newline normalization, since these arise from the publisher's packaging environment (e.g.,\ Windows) rather than the file content; we confirm with \texttt{diffoscope} that such residual differences are line-ending changes only.}


\section{Models for Independent Artifact Verifiability}
\label{sec:verifiability-model}

We model two aspects of independent artifact verification:
  a \textit{verifier model}, to describe the capabilities of an independent verifier,
  and
  an \textit{artifact comparison model}, describing how to compare published vs. rebuilt artifacts.

\JD{Three design choices: (1) How you recover information, and (2) How you build, and (3) How you compare. Maybe a good opportunity for a checks-and-exes or otherwise comparison table placed here in the Overview. Then A, B, and C of this section cover each topic.}



% An independent verifier can choose between two rebuild strategies that trade off security, automation, and success. \emph{Canonical rebuild} uses the ecosystem's standard build and packaging commands (\S\ref{sec:background}). It avoids
% replaying the original CI/CD workflow, reducing exposure to compromised workflow dependencies and helping surface artifacts that diverge from the expected source, but it can fail when a package relies on custom workflow steps or package-specific build logic. \emph{Workflow replay} instead re-executes the package's public CI/CD workflow, more closely approximating the maintainer's environment, but it inherits the workflow's risks and costs: compromised actions, irrelevant steps, long runtimes, and failures from missing secrets, tokens, or caches, and, as \S\ref{sec:rq2} shows, rarely succeeds in practice. We study how both choices
% affect verifiability.


% An independent verifier can choose between two rebuild strategies that trade off security, automation, and verification success.
% The first strategy uses ecosystem-level canonical build and packaging commands, such as \texttt{python -m build} for PyPI or \texttt{npm pack} for NPM.
% This strategy avoids replaying the original CI/CD workflow, reducing exposure to compromised workflow dependencies and helping detect artifacts that diverge from the expected source.
% However, it may fail when packages rely on custom workflow steps, undocumented setup, or package-specific build logic.

% The second strategy replays public CI/CD workflows, such as GitHub Actions.
% Workflow replay may improve verification success because it better approximates\footnote{Still an approximation, because workflows may interact with private machines whose configuration is unknown.} the maintainer's original build environment and commands.
% However, it also inherits the workflow's risks and costs: compromised actions, irrelevant steps, long runtimes, and failures caused by missing secrets, tokens, caches, or deployment credentials.
% \JD{next sentence is false right}
% In this work, we study how these rebuild choices affect independent artifact verifiability.
% \KC{should we have a figure to articulate this section/workflow/entity/feature definition of our model and its unique attributes?}
% \JD{Consider stronger criticism of prior work w.r.t. informality. If we are claiming a novel model (contribution bullet 1) then this next sentence is giving away that claim. If the prior work doesn't articulate cleanly then it's not a model}
% \JD{LastPyMile? Other things? This section model needs to explain how what we're proposin copmares to those prior famous SE papers}

\subsection{Verifier Model}
\label{subsec:verifier-model}

We model a \textbf{verifier} as an independent entity that validates a published artifact using only registry-exposed or registry-derived information, without access to the original build machine, private CI/CD state, or undisclosed secrets. Examples include registries performing publication-time checks and downstream users validating artifacts before installation.

The verifier makes two design choices: how to recover rebuild metadata and how to rebuild the artifact. 
Our model focuses on fully-automatable and generalizable metadata recovery.
It recovers metadata from three complementary sources: the published artifact, the registry-linked source repository, and registry-linked provenance when available. 
These sources provide evidence such as embedded source information, build-tool lock files, tags, commit history, manifests, source commits, and CI/CD workflows. 
We restrict recovery to these sources because they are available to independent verifiers and can scale across ecosystems and packages.

For rebuilding, our model adopts an independent builder approach~\cite{lamb_reproducible_2022}. 
Instead of replaying the original CI/CD workflow, the verifier uses ecosystem-level canonical build and packaging commands, such as \texttt{npm pack} for npm and \texttt{python -m build} for PyPI. 
This reduces dependence on potentially compromised workflow steps and tests whether artifacts can be reproduced from durable source and build metadata. 
However, it can fail when packages rely on custom workflow steps or package-specific build logic.

Unlike prior systems tailored to specific ecosystems or package sets, such as BuildGen for Maven~\cite{hassanshahiUnlockingReproducibilityAutomating2025} and OSS-Rebuild~\cite{google_oss_rebuild}, our model intentionally restricts the verifier to ecosystem-level, registry-derived evidence to measure how far independent artifact verification can scale across decentralized-build package ecosystems.


\subsection{Artifact Comparison Model}
\label{subsec:artifact-comparison-model}

Artifact verification requires a criterion for comparing the published artifact with the independently rebuilt artifact. 
Package artifacts in decentralized-build ecosystems are mostly structured archives with inspectable contents, rather than flat byte sequences as modeled by prior work~\cite{dietrich_levels_2025}.
As a result, we define an \textbf{artifact comparison model} that uses tiered equivalence levels over package contents and metadata.

\begin{itemize}
\item \textit{L1: byte-identical.} The published and rebuilt artifacts are byte-for-byte identical. This is the strongest level and matches the standard definition of reproducibility.

\item \textit{L2: content-identical.} The artifacts contain the same files with identical contents after normalizing archive-level metadata, such as timestamps, file ordering, compression parameters, ownership fields, and permissions.

\item \textit{L3: code-identical.} The artifacts match after excluding or normalizing known non-executable content differences, such as documentation, generated package metadata, runtime-irrelevant lock files, or non-executable fields in build definition files.

\item \textit{L4: behavioral equivalence.} The artifacts differ syntactically, but source-code, binary, or semantic analysis shows that the differences do not affect execution. We treat L4 as out of scope because such analyses may be unavailable, incomplete, or unsound for arbitrary package artifacts.
\end{itemize}

Unlike prior systems that report only L1 (bit-identical)~\cite{google_oss_rebuild} or rebuild success~\cite{hassanshahiUnlockingReproducibilityAutomating2025}, this model makes artifact comparison explicit and reports outcomes for multiple equivalence levels.
